% 第 14 章 Using a Milkyway Database
\bichapter{使用 Milkyway 数据库}{Using a Milkyway Database}
\label{chap:mw}

可在 DC Explorer 内写出 Milkyway 数据库，以便与其他基于 Synopsys Milkyway 数据库的工具（如 IC Compiler）配合使用。为此使用 \cmd{write\_milkyway} 命令。

使用 Milkyway 数据库时，无需使用 Verilog 或 VHDL 等中间网表文件交换格式与其他基于 Synopsys Milkyway 数据库的工具通信。
在 DC Explorer 内使用 Milkyway 数据库之前，必须准备设计库与参考库，并了解下列概念与任务：
\begin{itemize}
  \item 关于 Milkyway 数据库
  \item 使用 Milkyway 数据库的指南
  \item 创建 Milkyway 设计库
  \item 写出 Milkyway 数据库
\end{itemize}

% ============================================================
\section{关于 Milkyway 数据库}
\subsection*{About the Milkyway Database}

Milkyway 数据库将设计数据存储在 Milkyway 设计库中，将物理库数据存储在 Milkyway 参考库中。
\begin{itemize}
  \item \textbf{Milkyway 设计库}\\
        用于存储设计数据（即已 uniquify 的已映射网表与约束）的 Milkyway 目录结构称为 Milkyway 设计库。
        通过将变量 \varname{mw\_design\_library} 设为根目录路径，为当前会话指定 Milkyway 设计库。
  \item \textbf{Milkyway 参考库}\\
        用于存储物理库数据的 Milkyway 目录结构称为 Milkyway 参考库。
        参考库包含标准单元、宏单元与 pad 单元。关于创建参考库的信息，见 Milkyway 文档。\\
        通过将变量 \varname{mw\_reference\_library} 设为根目录路径，为当前会话指定 Milkyway 参考库。
        列表中的顺序表示参考冲突解析的优先级；若多个参考库有同名单元，则第一个参考库优先。
\end{itemize}

% ------------------------------------------------------------
\subsection{所需许可证与文件}
\subsubsection*{Required License and Files}

使用 Milkyway 数据库之前，需要下列所需许可证与文件：
\begin{itemize}
  \item Milkyway-Interface 许可证\\
        DC Explorer 提供此许可证，用于运行 \cmd{write\_milkyway} 命令。
  \item 逻辑库源（\texttt{.lib}）
  \item 已编译数据库
  \begin{itemize}
    \item 逻辑库（\texttt{.db}），包含标准单元时序、功耗、功能、测试等
    \item Milkyway 库（FRAM），包含技术数据
  \end{itemize}
\end{itemize}

% ------------------------------------------------------------
\subsection{调用 Milkyway Environment 工具}
\subsubsection*{Invoking the Milkyway Environment Tool}

要调用 Milkyway Environment 工具，输入：
\begin{lstlisting}
% Milkyway -galaxy
\end{lstlisting}

该命令检出 Milkyway-Interface 许可证。Milkyway Environment 工具是图形用户界面（GUI），用于操作 Milkyway 库。
可用该工具维护 Milkyway 设计库，例如删除设计的未使用版本。

关于使用 Milkyway Environment 工具的信息，见 Milkyway 文档。

\begin{seeAlsoBox}
\begin{itemize}
  \item 使用 Milkyway 数据库的指南
  \item 使用 Milkyway 格式的限制
  \item 写出 Milkyway 数据库
\end{itemize}
\end{seeAlsoBox}

% ============================================================
\section{使用 Milkyway 数据库的指南}
\subsection*{Guidelines for Using the Milkyway Databases}

使用 \cmd{write\_milkyway} 命令时，请遵循下列指南。
\begin{itemize}
  \item 确保 Milkyway 参考库中存在的所有单元在时序库中都有对应单元。
        Milkyway 参考库中单元的端口方向由时序库中单元的端口方向设置。
        若单元存在于 Milkyway 参考库但不在时序库中，则参考库中这些单元的端口方向不会被设置。
  \item 在运行 \cmd{write\_milkyway} 之前运行 \cmd{uniquify} 命令。
  \item 必须确保逻辑库与 Milkyway technology 文件中的单位一致。\\
        SDC 文件不含单位信息。若逻辑库与 Milkyway technology 文件中的单位不一致，\cmd{write\_milkyway} 无法自动转换。
        例如，若逻辑库使用飞法（femtofarad）作为电容单位，而 Milkyway technology 文件使用皮法（picofarad），则 \cmd{write\_sdc} 的输出会显示不同的网负载值。\\
        在下列示例中，逻辑库与 Milkyway technology 文件中的电容单位不一致。运行 \cmd{write\_milkyway} 之前，网 \texttt{gpdhi\_word\_d\_21\_} 显示下列 \cmd{set\_load} 信息：
\begin{lstlisting}
set_load 1425.15 [get_nets {gpdhi_word_d_21_}]
\end{lstlisting}
        运行 \cmd{write\_milkyway} 之后，SDC 文件显示：
\begin{lstlisting}
set_load 8.36909 [get_nets {gpdhi_word_d_21_}]
\end{lstlisting}
  \item DC Explorer 区分大小写。可在运行 \cmd{write\_milkyway} 之前通过下列任务以不区分大小写模式使用工具：
  \begin{itemize}
    \item 准备 link library 中所用库的大写版本。
    \item 使用 \cmd{change\_names} 命令确保网表为大写。
  \end{itemize}
\end{itemize}

% ============================================================
\section{创建 Milkyway 设计库}
\subsection*{Creating a Milkyway Design Library}

使用 Milkyway 设计库指定物理库并以 Milkyway 格式保存设计。
创建 Milkyway 设计库所需的输入为：Milkyway 参考库与 Milkyway technology 文件。

要创建 Milkyway 设计库：
\begin{enumerate}
  \item 用 \cmd{create\_mw\_lib} 命令创建 Milkyway 设计库，如下例所示：
\begin{lstlisting}
de_shell> create_mw_lib -technology $mw_tech_file \
   -mw_reference_library $mw_reference_library $mw_design_library_name
\end{lstlisting}
  \item 用 \cmd{open\_mw\_lib} 命令打开已创建的 Milkyway 库，如下例所示：
\begin{lstlisting}
de_shell> open_mw_lib $mw_design_library_name
\end{lstlisting}
  \item （可选）用 \cmd{set\_tlu\_plus\_files} 命令附着 TLUPlus 文件，如下例所示：
\begin{lstlisting}
de_shell> set_tlu_plus_files -max_tluplus $max_tlu_file \
   -min_tluplus $min_tlu_file -tech2itf_map $prs_map_file
\end{lstlisting}
  \item 在后续 \cmd{de\_shell} 会话中用 \cmd{open\_mw\_lib} 打开 Milkyway 库。
        若使用 TLUPlus 文件进行 RC 估算，用 \cmd{set\_tlu\_plus\_files} 附着这些文件，如下例所示：
\begin{lstlisting}
de_shell> open_mw_lib $mw_design_library_name
de_shell> set_tlu_plus_files -max_tluplus $max_tlu_file \
   -min_tluplus $min_tlu_file -tech2itf_map $prs_map_file
\end{lstlisting}
\end{enumerate}

\begin{seeAlsoBox}
关于 Milkyway 数据库
\end{seeAlsoBox}

% ============================================================
\section{写出 Milkyway 数据库}
\subsection*{Writing the Milkyway Database}

要将设计数据保存在 Milkyway 设计库中，使用 \cmd{write\_milkyway} 命令。
该命令将内存中的网表与 floorplan 信息写入 Milkyway 设计库格式，重新创建层次保留信息，并将当前设计的设计数据保存在设计文件中。
设计文件路径为 \texttt{design\_dir/CEL/file\_name:version}，其中 \texttt{design\_dir} 是变量 \varname{mw\_design\_library} 中指定的位置。

例如，用 \opt{-output} 选项指定文件名：
\begin{lstlisting}
de_shell> write_milkyway -output my_file -overwrite
\end{lstlisting}

要将内存中的设计信息写入名为 \texttt{testmw} 的 Milkyway 库并将设计文件命名为 \texttt{TOP}，输入：
\begin{lstlisting}
de_shell> set_app_var mw_design_library testmw
de_shell> write_milkyway -output TOP
\end{lstlisting}

示例~\ref{ex:mw-script} 给出设置并写出 Milkyway 数据库的脚本。

\begin{lstlisting}[caption={设置并写出 Milkyway 数据库的脚本},label=ex:mw-script]
set_app_var search_path "$search_path ./libraries"
set_app_var link_library "* max_lib.db"
set_app_var target_library "max_lib.db"

create_mw_lib -technology $mw_tech_file -mw_reference_library \
   $mw_reference_library$mw_lib_name
open_mw_lib $mw_lib_name
read_file -format ddc design.ddc
current_design TopDesign
link
write_milkyway -output myTop
\end{lstlisting}

\begin{seeAlsoBox}
\begin{itemize}
  \item 关于 \cmd{write\_milkyway} 命令
  \item 使用 Milkyway 数据库的指南
  \item 使用 Milkyway 格式的限制
  \item 使用层次模型
\end{itemize}
\end{seeAlsoBox}

% ------------------------------------------------------------
\subsection{关于 write\_milkyway 命令}
\subsubsection*{About the write\_milkyway Command}

使用 \cmd{write\_milkyway} 命令时，请记住下列要点：
\begin{itemize}
  \item 必须在运行 \cmd{write\_milkyway} 之前运行 \cmd{create\_mw\_lib} 命令。
  \item 若设计文件已存在（即已对同一输出目录的设计运行过 \cmd{write\_milkyway}），该命令会创建额外的设计文件并递增版本号。
        必须确保在 Milkyway 中打开正确版本；默认情况下 Milkyway 打开最新版本。
        为避免创建额外版本，指定 \opt{-overwrite} 选项以覆盖设计文件的当前版本并节省磁盘空间。
  \item 该命令不修改内存中的数据。
  \item 内存中设计上存在的、在 Milkyway 中有等价属性的属性会被转换（并非设计数据库中的所有属性都会被转换）。
  \item 用 \cmd{write\_milkyway} 转换的层次网表在其 Milkyway 数据库中保留层次。
\end{itemize}

\begin{seeAlsoBox}
写出 Milkyway 数据库
\end{seeAlsoBox}

% ------------------------------------------------------------
\subsection{使用 Milkyway 格式的限制}
\subsubsection*{Limitations of Using Milkyway Format}

以 Milkyway 格式写出设计时适用下列限制：
\begin{itemize}
  \item 设计必须已映射。\\
        由于 Milkyway 格式描述物理信息，它仅支持已映射设计。不能用 Milkyway 格式存储未映射设计的设计数据。
  \item 设计不得包含多重实例。\\
        必须以 Milkyway 格式保存之前对设计做 uniquify。使用 \cmd{check\_design -multiple\_designs} 命令报告与多重实例化设计相关的信息。
  \item \cmd{write\_milkyway} 将整个层次设计保存在单个 Milkyway 设计文件中。不能为每个子设计生成单独的设计文件。
\end{itemize}

\begin{seeAlsoBox}
\begin{itemize}
  \item 关于 Milkyway 数据库
  \item 写出 Milkyway 数据库
\end{itemize}
\end{seeAlsoBox}
